% Cronología de las campañas del puente de Palma del Río que hay en el repositorio.
% Círculo lleno y azul: campaña con registros de aceleración. Círculo hueco y gris: solo fotos.
\documentclass[tikz,border=4pt]{standalone}
\usepackage{resonancia}
\begin{document}
\begin{tikzpicture}
  % eje del tiempo, con un corte entre 2014 y 2020 (la escala no es proporcional)
  \draw[eje] (0,0) -- (3.0,0);
  \draw[eje] (3.5,0) -- (12.6,0);
  \draw[trazo, color=rGris] (3.05,-0.18) -- (3.15,0.18);
  \draw[trazo, color=rGris] (3.35,-0.18) -- (3.45,0.18);
  \node[etiqueta, color=rGris, anchor=north] at (3.25,-0.25) {\scriptsize corte de escala};

  % 2014-06: solo fotos
  \draw[trazo, color=rGris, fill=white] (1.4,0) circle (0.14);
  \node[etiqueta, anchor=south] at (1.4,0.25) {\textbf{jun. 2014}};
  \node[etiqueta, anchor=south, align=center, color=rGris] at (1.4,0.75) {campaña modal\\4 fotos};

  % 2020-07-17: datos
  \fill[rAzul] (5.0,0) circle (0.14);
  \node[etiqueta, anchor=south] at (5.0,0.25) {\textbf{17 jul. 2020}};
  \node[etiqueta, anchor=south, align=center, color=rAzul] at (5.0,0.75) {operacional\\9 registros\\512 Hz};

  % 2020-10-16: datos
  \fill[rAzul] (8.0,0) circle (0.14);
  \node[etiqueta, anchor=south] at (8.0,0.25) {\textbf{16 oct. 2020}};
  \node[etiqueta, anchor=south, align=center, color=rAzul] at (8.0,0.75) {operacional\\3 registros\\2048 Hz};

  % 2021-04-18: solo fotos
  \draw[trazo, color=rGris, fill=white] (11.0,0) circle (0.14);
  \node[etiqueta, anchor=south] at (11.0,0.25) {\textbf{18 abr. 2021}};
  \node[etiqueta, anchor=south, align=center, color=rGris] at (11.0,0.75) {vistas\\3 fotos};

  % leyenda
  \fill[rAzul] (0.3,-0.95) circle (0.11);
  \node[etiqueta, anchor=west] at (0.5,-0.95) {con registros de aceleración};
  \draw[trazo, color=rGris, fill=white] (5.6,-0.95) circle (0.11);
  \node[etiqueta, anchor=west] at (5.8,-0.95) {solo fotos};
\end{tikzpicture}
\end{document}
